\documentclass[aspectratio=169]{beamer}

% ==== Theme & basic look ====
\usetheme{Madrid}
\usecolortheme{seahorse}
\setbeamertemplate{navigation symbols}{}
\setbeamertemplate{footline}[frame number]


% Packages for mathematics  
\usepackage{amsmath}  
\usepackage{amsfonts}  
\usepackage{amssymb}  
\usepackage{CJKutf8}
% Support for \Tilde macro used throughout
\providecommand{\Tilde}[1]{\tilde{#1}}
\providecommand{\proofname}{Proof}
\newenvironment<>{proofs}[1][\proofname]{%
    \par
    \def\insertproofname{#1.}%
    \usebeamertemplate{proof begin}#2}
  {\usebeamertemplate{proof end}}
\makeatother

% Math shorthand and vector symbols
\newcommand{\grad}{\nabla}
\newcommand{\EE}{\mathbb{E}}
\newcommand{\dd}{\mathrm{d}}
\newcommand{\kbT}{k_B T}
\newcommand{\vx}{\mathbf{x}}
\newcommand{\vz}{\mathbf{z}}
\newcommand{\KL}{\text{KL}}
\newcommand{\ELBO}{\text{ELBO}}
\DeclareMathOperator{\sg}{sg}

% Bold vector/time-indexed variables
\newcommand{\vxz}{\mathbf{x}}
\newcommand{\vzt}{\mathbf{z}}
\newcommand{\vtheta}{\boldsymbol{\theta}}
\newcommand{\vphi}{\boldsymbol{\phi}}
\newcommand{\veps}{\boldsymbol{\epsilon}}

% Title page details:   
\title[Understanding VAE]{Understanding Variational Autoencoders and their Variants}  
\author{Pipi Hu}  
\institute[BIMSA]{Beijing Institute of Mathematical Sciences and Applications}  
\date{\today}  

\pgfdeclareimage[width=2cm]{logo}{figs/logo.png}
\logo{\pgfuseimage{logo}}

% Outline slide (moved to later; avoid frame before \begin{document})

\setbeamertemplate{navigation symbols}{}
\AtBeginSection[]  
{  
    \begin{frame}<beamer>{Outline}         
        \tableofcontents[currentsection]  
    \end{frame}  
}  
  
\AtBeginSubsection[]  
{  
    \begin{frame}<beamer>{Outline}         
        \tableofcontents[currentsection,currentsubsection]  
    \end{frame}  
}  

\begin{document}  
  
% Title slide  
\begin{frame}  
  \titlepage  
\end{frame}  

\begin{frame}{Feynman's words}
\centering
    What I cannot create, I do not understand.
\end{frame}
  
% Presentation slides  

\section{Introduction: The Problem of Generative Modeling}

\begin{frame}{The Challenge of Generative Modeling}
Given a dataset $\mathcal{D} = \{x^{(i)}\}_{i=1}^N$ sampled from an unknown distribution $p^*(x)$, how can we:
\begin{itemize}
    \item Learn a model $p_\theta(x)$ that approximates $p^*(x)$?
    \item Generate new samples from the learned distribution?
    \item Perform inference: $p_\theta(z|x)$ for latent variables $z$?
\end{itemize}

\textcolor{red}{Key challenges:}
\begin{enumerate}
    \item High-dimensional data spaces
    \item Intractable likelihood computation
    \item Latent variable inference
    \item Scalable training procedures
\end{enumerate}
\end{frame}

\begin{frame}{Traditional Approaches and Their Limitations}
\begin{itemize}
    \item \textbf{Maximum Likelihood Estimation (MLE):}
    \begin{equation}
        \theta^* = \arg\max_\theta \sum_{i=1}^N \log p_\theta(x^{(i)})
    \end{equation}
    \item \textbf{Problem:} Direct likelihood computation often intractable
    \item \textbf{Latent Variable Models:}
    \begin{equation}
        p_\theta(x) = \int p_\theta(x|z) p_\theta(z) dz
    \end{equation}
    \item \textbf{Problem:} Integration over latent space $z$ is typically intractable
\end{itemize}

\textcolor{red}{Solution:} Variational inference to approximate intractable posteriors
\end{frame}

\section{Variational Inference: The Foundation}

\begin{frame}{Variational Inference Framework}
Given a latent variable model $p_\theta(x,z) = p_\theta(x|z)p_\theta(z)$, we want to:
\begin{itemize}
    \item Approximate the intractable posterior $p_\theta(z|x)$
    \item Maximize the marginal likelihood $p_\theta(x)$
\end{itemize}

\textbf{Key idea:} Introduce a variational distribution $q_\phi(z|x)$ to approximate $p_\theta(z|x)$

\begin{equation}
\KL(q_\phi(z|x) \| p_\theta(z|x)) = \EE_{q_\phi(z|x)}[\log q_\phi(z|x)] - \EE_{q_\phi(z|x)}[\log p_\theta(z|x)]
\end{equation}
\end{frame}

\begin{frame}{Evidence Lower Bound (ELBO)}
The KL divergence can be decomposed as:
\begin{equation}
\KL(q_\phi(z|x) \| p_\theta(z|x)) = \log p_\theta(x) - \ELBO(x; \theta, \phi)
\end{equation}

where the \textcolor{red}{Evidence Lower Bound} is:
\begin{equation}
\ELBO(x; \theta, \phi) = \EE_{q_\phi(z|x)}[\log p_\theta(x|z)] - \KL(q_\phi(z|x) \| p_\theta(z))
\end{equation}

\textbf{Key insight:} Since $\KL \geq 0$, we have:
\begin{equation}
\log p_\theta(x) \geq \ELBO(x; \theta, \phi)
\end{equation}

Maximizing ELBO $\Leftrightarrow$ Minimizing KL divergence $\Leftrightarrow$ Approximating posterior
\end{frame}

\begin{frame}{ELBO Interpretation}
The ELBO can be rewritten as:
\begin{equation}
\ELBO(x; \theta, \phi) = \EE_{q_\phi(z|x)}[\log p_\theta(x|z)] - \KL(q_\phi(z|x) \| p_\theta(z))
\end{equation}

\textbf{Two components:}
\begin{enumerate}
    \item \textcolor{blue}{Reconstruction term:} $\EE_{q_\phi(z|x)}[\log p_\theta(x|z)]$
    \begin{itemize}
        \item Encourages $q_\phi(z|x)$ to encode information about $x$
        \item Maximizes likelihood of reconstructing $x$ from $z$
    \end{itemize}
    \item \textcolor{red}{Regularization term:} $-\KL(q_\phi(z|x) \| p_\theta(z))$
    \begin{itemize}
        \item Encourages $q_\phi(z|x)$ to be close to prior $p_\theta(z)$
        \item Prevents overfitting to training data
    \end{itemize}
\end{enumerate}
\end{frame}

\section{Variational Autoencoders: Theory}

\begin{frame}{VAE Architecture}
\begin{center}
\textbf{Encoder-Decoder Architecture}
\end{center}

\begin{itemize}
    \item \textbf{Encoder (Recognition model):} $q_\phi(z|x)$
    \begin{itemize}
        \item Maps data $x$ to latent representation $z$
        \item Typically: $q_\phi(z|x) = \mathcal{N}(z; \mu_\phi(x), \sigma_\phi^2(x)I)$
    \end{itemize}
    \item \textbf{Decoder (Generative model):} $p_\theta(x|z)$
    \begin{itemize}
        \item Reconstructs data $x$ from latent $z$
        \item Typically: $p_\theta(x|z) = \mathcal{N}(x; \mu_\theta(z), \sigma_\theta^2 I)$ or Bernoulli
    \end{itemize}
    \item \textbf{Prior:} $p_\theta(z) = \mathcal{N}(z; 0, I)$
\end{itemize}
\end{frame}

\begin{frame}{The Reparameterization Trick}
\textbf{Problem:} How to backpropagate through stochastic sampling $z \sim q_\phi(z|x)$?

\textbf{Solution:} Reparameterization trick
\begin{equation}
z = \mu_\phi(x) + \sigma_\phi(x) \odot \epsilon, \quad \epsilon \sim \mathcal{N}(0, I)
\end{equation}

\textbf{Benefits:}
\begin{itemize}
    \item Gradient flows through deterministic path: $x \to \mu_\phi(x), \sigma_\phi(x) \to z$
    \item Stochasticity moved to $\epsilon$ (independent of parameters)
    \item Enables end-to-end training with backpropagation
\end{itemize}

\textbf{Alternative approaches:}
\begin{itemize}
    \item Score function estimator (higher variance)
    \item Control variates
\end{itemize}
\end{frame}

\begin{frame}{VAE Training Objective}
For a dataset $\mathcal{D} = \{x^{(i)}\}_{i=1}^N$, the VAE loss is:
\begin{equation}
\mathcal{L}(\theta, \phi) = \sum_{i=1}^N \ELBO(x^{(i)}; \theta, \phi)
\end{equation}

Expanding the ELBO:
\begin{equation}
\mathcal{L}(\theta, \phi) = \sum_{i=1}^N \left[ \EE_{q_\phi(z|x^{(i)})}[\log p_\theta(x^{(i)}|z)] - \KL(q_\phi(z|x^{(i)}) \| p_\theta(z)) \right]
\end{equation}

\textbf{Training procedure:}
\begin{enumerate}
    \item Sample $x^{(i)} \sim \mathcal{D}$
    \item Sample $\epsilon \sim \mathcal{N}(0, I)$
    \item Compute $z^{(i)} = \mu_\phi(x^{(i)}) + \sigma_\phi(x^{(i)}) \odot \epsilon$
    \item Compute reconstruction loss: $\log p_\theta(x^{(i)}|z^{(i)})$
    \item Compute KL divergence: $\KL(q_\phi(z|x^{(i)}) \| p_\theta(z))$
    \item Backpropagate gradients
\end{enumerate}
\end{frame}

\begin{frame}{Analytical KL Divergence}
For Gaussian distributions, the KL divergence has a closed form:
\begin{equation}
\KL(\mathcal{N}(\mu, \sigma^2) \| \mathcal{N}(0, I)) = \frac{1}{2} \sum_{j=1}^J (\mu_j^2 + \sigma_j^2 - \log \sigma_j^2 - 1)
\end{equation}

where $J$ is the dimension of the latent space.

\textbf{Interpretation:}
\begin{itemize}
    \item $\mu_j^2$: Penalizes large mean values
    \item $\sigma_j^2 - \log \sigma_j^2 - 1$: Encourages $\sigma_j \approx 1$
    \item Total effect: Regularizes latent space to be close to standard Gaussian
\end{itemize}

\textbf{Benefits:}
\begin{itemize}
    \item No sampling required for KL term
    \item Stable gradients
    \item Interpretable regularization
\end{itemize}
\end{frame}

\section{VAE Variants and Extensions}

\begin{frame}{$\beta$-VAE: Controlling the Regularization}
\textbf{Motivation:} Balance between reconstruction and regularization

\textbf{$\beta$-VAE objective:}
\begin{equation}
\mathcal{L}(\theta, \phi) = \EE_{q_\phi(z|x)}[\log p_\theta(x|z)] - \beta \KL(q_\phi(z|x) \| p_\theta(z))
\end{equation}

\textbf{Effect of $\beta$:}
\begin{itemize}
    \item $\beta > 1$: Stronger regularization $\Rightarrow$ Better disentanglement
    \item $\beta < 1$: Weaker regularization $\Rightarrow$ Better reconstruction
    \item $\beta = 1$: Standard VAE
\end{itemize}

\textbf{Disentanglement:} $\beta$-VAE encourages learning of independent latent factors
\end{frame}

\begin{frame}{VQ-VAE: Vector Quantized VAE}
\textbf{Key idea:} Use discrete latent representations instead of continuous

\textbf{Architecture:}
\begin{itemize}
    \item Encoder: $z_e = E(x)$ (continuous)
    \item Quantization: $z_q = \text{Quantize}(z_e)$ (discrete)
    \item Decoder: $\hat{x} = D(z_q)$
\end{itemize}

\textbf{Quantization process:}
\begin{equation}
z_q = \arg\min_{k} \|z_e - e_k\|^2
\end{equation}
where $\{e_k\}_{k=1}^K$ are learnable codebook vectors.

\textbf{Benefits:}
\begin{itemize}
    \item Discrete representations (useful for language, etc.)
    \item Better compression
    \item Enables autoregressive modeling in latent space
\end{itemize}
\end{frame}

\begin{frame}{VQ-VAE Training}
\textbf{Challenge:} Quantization is non-differentiable

\textbf{Solution:} Straight-through estimator
\begin{equation}
\frac{\partial \mathcal{L}}{\partial z_e} = \frac{\partial \mathcal{L}}{\partial z_q}
\end{equation}

\textbf{Training objective:}
\begin{equation}
\mathcal{L} = \log p(x|z_q) + \|\sg(z_e) - e_k\|^2 + \beta\|z_e - \sg(e_k)\|^2
\end{equation}

where $\sg(\cdot)$ is stop-gradient operator.

\textbf{Components:}
\begin{itemize}
    \item Reconstruction loss: $\log p(x|z_q)$
    \item Codebook loss: $\|\sg(z_e) - e_k\|^2$ (updates codebook)
    \item Commitment loss: $\beta\|z_e - \sg(e_k)\|^2$ (encourages encoder to commit)
\end{itemize}
\end{frame}

\begin{frame}{WAE: Wasserstein Autoencoder}
\textbf{Motivation:} Replace KL divergence with Wasserstein distance

\textbf{WAE objective:}
\begin{equation}
\mathcal{L} = \EE_{p(x)} \EE_{q_\phi(z|x)} [c(x, D(z))] + \lambda \mathcal{D}(q_\phi(z), p(z))
\end{equation}

where:
\begin{itemize}
    \item $c(x, D(z))$: Reconstruction cost
    \item $\mathcal{D}(q_\phi(z), p(z))$: Distance between aggregated posterior and prior
    \item $\lambda$: Regularization weight
\end{itemize}

\textbf{Benefits:}
\begin{itemize}
    \item More flexible than KL divergence
    \item Better mode coverage
    \item Can use different distance metrics
\end{itemize}
\end{frame}

\begin{frame}{VAE-GAN: Combining VAE and GAN}
\textbf{Idea:} Use GAN discriminator to improve VAE reconstruction quality

\textbf{Architecture:}
\begin{itemize}
    \item VAE encoder-decoder: $x \to z \to \hat{x}$
    \item GAN discriminator: Distinguishes real vs reconstructed samples
    \item Generator: VAE decoder
\end{itemize}

\textbf{Training objective:}
\begin{equation}
\mathcal{L} = \mathcal{L}_{\text{VAE}} + \lambda \mathcal{L}_{\text{GAN}}
\end{equation}

where:
\begin{itemize}
    \item $\mathcal{L}_{\text{VAE}}$: Standard VAE loss
    \item $\mathcal{L}_{\text{GAN}}$: Adversarial loss for discriminator
\end{itemize}

\textbf{Benefits:}
\begin{itemize}
    \item Better sample quality than VAE alone
    \item Maintains inference capabilities
\end{itemize}
\end{frame}

\section{Advanced Topics}

\begin{frame}{Hierarchical VAE}
\textbf{Motivation:} Learn hierarchical latent representations

\textbf{Architecture:}
\begin{equation}
p_\theta(x, z_1, \ldots, z_L) = p_\theta(x|z_1) \prod_{l=1}^{L-1} p_\theta(z_l|z_{l+1}) p_\theta(z_L)
\end{equation}

\textbf{Inference:}
\begin{equation}
q_\phi(z_{1:L}|x) = \prod_{l=1}^L q_\phi(z_l|z_{<l}, x)
\end{equation}

\textbf{Benefits:}
\begin{itemize}
    \item Multi-scale representations
    \item Better modeling of complex data
    \item Hierarchical generation process
\end{itemize}
\end{frame}

\begin{frame}{Conditional VAE (C-VAE)}
\textbf{Extension:} Incorporate conditioning information

\textbf{Model:}
\begin{equation}
p_\theta(x, z|c) = p_\theta(x|z, c) p_\theta(z)
\end{equation}

\textbf{Inference:}
\begin{equation}
q_\phi(z|x, c) = \mathcal{N}(z; \mu_\phi(x, c), \sigma_\phi^2(x, c))
\end{equation}

\textbf{Applications:}
\begin{itemize}
    \item Conditional generation
    \item Controlled synthesis
    \item Multi-modal learning
\end{itemize}

\textbf{Training:}
\begin{equation}
\mathcal{L} = \EE_{q_\phi(z|x,c)}[\log p_\theta(x|z,c)] - \KL(q_\phi(z|x,c) \| p_\theta(z))
\end{equation}
\end{frame}

\begin{frame}{Importance Weighted Autoencoder (IWAE)}
\textbf{Motivation:} Better approximation of log-likelihood

\textbf{IWAE objective:}
\begin{equation}
\mathcal{L}_K = \EE_{z_1, \ldots, z_K \sim q_\phi(z|x)} \left[ \log \frac{1}{K} \sum_{k=1}^K \frac{p_\theta(x, z_k)}{q_\phi(z_k|x)} \right]
\end{equation}

\textbf{Key properties:}
\begin{itemize}
    \item $\mathcal{L}_K \geq \mathcal{L}_{K-1}$ (monotonic improvement)
    \item $\lim_{K \to \infty} \mathcal{L}_K = \log p_\theta(x)$ (exact likelihood)
    \item $K=1$ recovers standard VAE
\end{itemize}

\textbf{Trade-offs:}
\begin{itemize}
    \item Better approximation with larger $K$
    \item Higher computational cost
    \item More complex gradients
\end{itemize}
\end{frame}

\section{Implementation and Practical Considerations}

\begin{frame}{VAE Implementation Details}
\textbf{Network architectures:}
\begin{itemize}
    \item \textbf{Encoder:} $x \to \mu_\phi(x), \log \sigma_\phi^2(x)$
    \item \textbf{Decoder:} $z \to \mu_\theta(z)$ (for Gaussian) or logits (for Bernoulli)
    \item \textbf{Activation:} ReLU, ELU, or Swish
    \item \textbf{Normalization:} BatchNorm or LayerNorm
\end{itemize}

\textbf{Training tricks:}
\begin{itemize}
    \item \textbf{KL annealing:} Gradually increase KL weight
    \item \textbf{Free bits:} Prevent posterior collapse
    \item \textbf{Warm-up:} Start with reconstruction only
\end{itemize}
\end{frame}

\begin{frame}{Common Issues and Solutions}
\textbf{Posterior collapse:}
\begin{itemize}
    \item \textbf{Symptom:} $q_\phi(z|x) \approx p_\theta(z)$ (no information in latent)
    \item \textbf{Causes:} Strong decoder, weak encoder
    \item \textbf{Solutions:} KL annealing, free bits, stronger encoder
\end{itemize}

\textbf{Blurry reconstructions:}
\begin{itemize}
    \item \textbf{Cause:} Gaussian decoder assumption
    \item \textbf{Solutions:} VAE-GAN, better decoder architectures
\end{itemize}

\textbf{Mode collapse:}
\begin{itemize}
    \item \textbf{Symptom:} Limited diversity in generated samples
    \item \textbf{Solutions:} Better regularization, architectural improvements
\end{itemize}
\end{frame}

\begin{frame}{Evaluation Metrics}
\textbf{Reconstruction quality:}
\begin{itemize}
    \item MSE, SSIM, FID (Fréchet Inception Distance)
    \item Perceptual losses
\end{itemize}

\textbf{Generation quality:}
\begin{itemize}
    \item FID, IS (Inception Score)
    \item Human evaluation
\end{itemize}

\textbf{Representation quality:}
\begin{itemize}
    \item Disentanglement metrics ($\beta$-VAE metric, MIG)
    \item Downstream task performance
    \item Interpolation smoothness
\end{itemize}
\end{frame}

\section{Recent Advances and Future Directions}

\begin{frame}{Recent Advances}
\textbf{Transformer-based VAEs:}
\begin{itemize}
    \item VQ-VAE-2: Hierarchical discrete representations
    \item DALL-E: Text-to-image generation
    \item VQGAN: High-quality image generation
\end{itemize}

\textbf{Improved training:}
\begin{itemize}
    \item NVAE: Normalizing flow decoder
    \item VDVAE: Very deep VAE
    \item VDM: Diffusion-based decoder
\end{itemize}

\textbf{Applications:}
\begin{itemize}
    \item Text generation (GPT-style models)
    \item Image synthesis and editing
    \item Audio generation
    \item Scientific modeling
\end{itemize}
\end{frame}

\begin{frame}{Open Questions and Future Directions}
\textbf{Theoretical understanding:}
\begin{itemize}
    \item Why do VAEs work despite approximation errors?
    \item Optimal choice of prior distributions
    \item Connection to other generative models
\end{itemize}

\textbf{Technical challenges:}
\begin{itemize}
    \item Better posterior approximations
    \item Scalable training for large models
    \item Integration with other architectures
\end{itemize}

\textbf{Applications:}
\begin{itemize}
    \item Multi-modal learning
    \item Scientific discovery
    \item Efficient compression
    \item Robust representations
\end{itemize}
\end{frame}

\section{Conclusion}

\begin{frame}{Summary}
\textbf{Key contributions of VAE:}
\begin{itemize}
    \item Scalable variational inference for latent variable models
    \item End-to-end training with reparameterization trick
    \item Foundation for many generative models
\end{itemize}

\textbf{VAE variants address:}
\begin{itemize}
    \item Disentanglement ($\beta$-VAE)
    \item Discrete representations (VQ-VAE)
    \item Sample quality (VAE-GAN, WAE)
    \item Hierarchical modeling (HVAE)
\end{itemize}

\textbf{Impact:}
\begin{itemize}
    \item Influenced modern generative models
    \item Enabled practical latent variable modeling
    \item Foundation for many applications
\end{itemize}
\end{frame}

\begin{frame}{Questions}
If you are interested in the questions, please feel free to write an email to me {\color{red}(\href{mailto:pisquare@microsoft.com}{pisquare@microsoft.com})} with your comments. 
\begin{enumerate}[(1.)]
    \item How does the choice of prior $p_\theta(z)$ affect VAE performance? What are the trade-offs between different priors?
    \item Can you derive the gradient of the ELBO with respect to the encoder parameters $\phi$? How does the reparameterization trick make this tractable?
    \item What are the theoretical guarantees of VAE? Under what conditions does the ELBO provide a good approximation to the true log-likelihood?
    \item How would you extend VAE to handle sequential data (like in language modeling)? What are the key challenges?
\end{enumerate}
\end{frame}

% Thank you slide  
\begin{frame}{Welcome inputs}  
  \centering  
  Any comments?
  \\[2em]
  Welcome your inputs!  
\end{frame}  
  
\end{document}
